\documentclass[11pt]{amsart}
\usepackage[T1]{fontenc}
\usepackage{amsmath,amssymb,amsthm}
\usepackage{graphicx}
\usepackage[colorlinks=true,linkcolor=blue,citecolor=blue,urlcolor=blue]{hyperref}

\newtheorem{theorem}{Theorem}
\newtheorem{lemma}{Lemma}

\newcommand{\K}{\mathbb K}
\newcommand{\Lop}{\mathcal L}
\newcommand{\QMA}{\operatorname{QMA}}
\newcommand{\MA}{\operatorname{MA}}
\newcommand{\Id}{\operatorname{Id}}
\newcommand{\vertiii}[1]{{\left\vert\kern-0.25ex\left\vert\kern-0.25ex\left\vert #1
\right\vert\kern-0.25ex\right\vert\kern-0.25ex\right\vert}}

\begin{document}

\title{A Fixed-Range Obstruction for Quasi Minimum-Attaining Operators}
\author{Run fa\_banach\_001}
\date{2026-06-16}
\maketitle

\section*{Status}

This is a candidate full answer to the open question in Section 2, page 9,
of Garcia--Maestre--Martin--Roldan \cite{GMMR}. The answer is positive:
for every infinite-dimensional Banach range space \(Y\), there is a Banach
space \(X\) such that
\[
  \QMA(X,Y)\neq \Lop(X,Y).
\]
In fact, \(X\) can be chosen to be an equivalent renorming of the same
underlying space \(Y\).

\section*{Source Question}

The source paper recalls that, for every infinite-dimensional domain space
\(X\), one can find a range space \(Y\) with
\(\QMA(X,Y)\neq\Lop(X,Y)\), and then asks for the reversed, fixed-range
form:

\begin{quote}
Is it true that for every infinite-dimensional Banach space \(Y\), there is
a Banach space \(X\) such that \(\QMA(X,Y)\neq\Lop(X,Y)\)?
\end{quote}

\begin{center}
\includegraphics[width=\textwidth]{figures/open_problem_crop.png}
\end{center}

\section*{Definitions}

For \(T\in\Lop(X,Y)\), its minimum norm is
\[
  m(T)=\inf\{\|Tx\|:x\in S_X\}.
\]
The operator \(T\) is quasi minimum-attaining if there are
\((x_k)\subset S_X\) and \(y\in m(T)S_Y\) such that \(Tx_k\to y\).
The class of such operators is denoted by \(\QMA(X,Y)\).

\section*{Proof Intuition}

The obstruction is the inverse of a non-norm-attaining identity map.  The
Josefson--Nissenzweig theorem gives a weak-star null sequence
\((y_n^*)\subset S_{Y^*}\).  Weighting these functionals by scalars
\(t_n\uparrow1\) defines an equivalent norm
\[
  \vertiii{y}=\max\left\{a\|y\|,\sup_n t_n |y_n^*(y)|\right\},
  \qquad 0<a<1.
\]
The identity from the original \(Y\) into \((Y,\vertiii{\cdot})\) has norm
one, because the weights approach one, but it cannot attain that norm:
for each fixed \(y\), weak-star nullity forces \(y_n^*(y)\to0\), while every
finite weight is strictly below one.  Inverting this identity gives an
isomorphism from the renormed space into the original \(Y\).  If that
isomorphism were QMA, the closed-range convergence would make the inverse
attain its norm, contradicting the construction.

\section*{The Result}

\begin{theorem}
Let \(Y\) be an infinite-dimensional real or complex Banach space. Then
there is an equivalent norm \(\vertiii{\cdot}\) on the vector space \(Y\)
such that, for \(X=(Y,\vertiii{\cdot})\), the identity operator
\[
  T=\Id:X\longrightarrow Y
\]
is bounded below but does not belong to \(\QMA(X,Y)\). Consequently, the
inclusion \(\QMA(X,Y)\subset \Lop(X,Y)\) is proper.
\end{theorem}

\begin{proof}
By the Josefson--Nissenzweig theorem, there is a weak-star null sequence
\((y_n^*)_{n\geq1}\subset S_{Y^*}\). Fix \(0<a<1\), and put
\(t_n=n/(n+1)\). Define, for \(y\in Y\),
\[
  \vertiii{y}
  =\max\left\{a\|y\|,\sup_{n\geq1} t_n |y_n^*(y)|\right\}.
\]
This is an equivalent norm on \(Y\), since
\[
  a\|y\|\leq \vertiii{y}\leq \|y\|
  \qquad (y\in Y).
\]

Let
\[
  J=\Id:Y\longrightarrow (Y,\vertiii{\cdot})
\]
denote the identity map from the original normed space into the renormed
space. We first show that \(\|J\|=1\) and that \(J\) does not attain its
norm.

If \(\|y\|\leq1\), then \(\vertiii{y}\leq1\), so \(\|J\|\leq1\). Conversely,
given \(n\) and \(\delta>0\), since \(\|y_n^*\|=1\), there is
\(y\in S_Y\) such that \(|y_n^*(y)|>1-\delta\). Hence
\[
  \|J y\|=\vertiii{y}\geq t_n(1-\delta).
\]
Letting \(n\to\infty\) and \(\delta\downarrow0\) gives \(\|J\|=1\).

Now fix \(y\in B_Y\). Since \(y_n^*(y)\to0\), choose \(N\geq1\) such that
\(|y_n^*(y)|\leq1/2\) for all \(n>N\). Then
\[
  \sup_{n>N} t_n|y_n^*(y)|\leq\frac12,
  \qquad
  \max_{1\leq n\leq N} t_n|y_n^*(y)|\leq t_N<1,
\]
while \(a\|y\|\leq a<1\). Therefore
\[
  \vertiii{y}
  \leq \max\left\{a,t_N,\frac12\right\}<1.
\]
Thus no point of \(B_Y\) attains the norm of \(J\).

Set \(X=(Y,\vertiii{\cdot})\), and let \(T=\Id:X\to Y\). Then \(T\) is an
isomorphism and \(T^{-1}=J\). Hence
\[
  m(T)=\frac{1}{\|T^{-1}\|}=1.
\]
The operator \(T\) cannot attain its minimum norm: if some \(x\in S_X\)
satisfied \(\|Tx\|=m(T)=1\), then \(Tx\in S_Y\) and
\[
  \|J(Tx)\|_X=\|x\|_X=1,
\]
so \(J\) would attain its norm, impossible.

It remains only to exclude quasi minimum attainment. Suppose, toward a
contradiction, that \(T\in\QMA(X,Y)\). Since \(m(T)=1\), there are
\((x_k)\subset S_X\) and \(y\in S_Y\) such that \(Tx_k\to y\) in \(Y\).
Applying the continuous inverse \(J=T^{-1}\), we get
\[
  x_k=J(Tx_k)\longrightarrow J y
  \quad\text{in } X.
\]
Therefore \(\|J y\|_X=\lim_k\|x_k\|_X=1\), with \(y\in S_Y\). This says
that \(J\) attains its norm, the contradiction. Consequently
\(T\notin\QMA(X,Y)\), proving the theorem.
\end{proof}

\section*{Verification Notes}

The proof is analytic. The only external theorem used is the
Josefson--Nissenzweig theorem, in the standard form that every
infinite-dimensional Banach space \(Y\) admits a weak-star null sequence in
\(S_{Y^*}\).

The last paragraph also gives a direct replacement for the usual closed-range
fact: for bounded-below operators, quasi minimum-attainment implies
minimum-attainment.

\section*{Novelty Check}

Local run indexes were searched for the arXiv id \(2405.01302\), the phrases
``quasi minimum-attaining'' and ``minimum-attaining'', the expression
``QMA of \(X,Y\)'', and ``Josefson--Nissenzweig''.
No duplicate packet or attempt for this question was found. External web
search on 2026-06-16 used the exact title of \cite{GMMR}, the arXiv id, and
phrases including ``quasi minimum-attaining for every infinite-dimensional
\(Y\)'', ``QMA(X,Y) minimum norm Banach Josefson'', and
``minimum-attaining Josefson-Nissenzweig quasi''. The search found the source
paper but no later paper explicitly answering this fixed-range question.

\section*{Limitations and Human Review Recommendation}

This solves the existential fixed-range question. It does not address the
separate implication questions around Lindenstrauss properties \(A\), \(A_M\),
and \(B_M\) mentioned earlier in the source paper.

For review, check the non-attainment argument for the renormed identity
\(J:Y\to X\), especially the use of weak-star nullity to make
\(\sup_n t_n|y_n^*(y)|<1\) for each fixed \(y\in B_Y\). Once that point is
accepted, the QMA contradiction is immediate.

\begin{thebibliography}{9}

\bibitem{GMMR}
D. Garcia, M. Maestre, M. Martin, and O. Roldan,
\emph{On density and Bishop-Phelps-Bollobas type properties for the minimum norm},
arXiv:2405.01302, 2024.
\url{https://arxiv.org/abs/2405.01302}.

\bibitem{Diestel}
J. Diestel,
\emph{Sequences and Series in Banach Spaces},
Graduate Texts in Mathematics, vol. 92, Springer, 1984.

\end{thebibliography}

\end{document}
